\documentclass{article}
\usepackage[final]{neurips_2025}
\usepackage[utf8]{inputenc}
\usepackage[T1]{fontenc}
\usepackage{hyperref}
\usepackage{url}
\usepackage{booktabs}
\usepackage{amsfonts}
\usepackage{amsmath}
\usepackage{amssymb}
\usepackage{amsthm}
\newtheorem{theorem}{Theorem}
\usepackage{nicefrac}
\usepackage{microtype}
\usepackage{graphicx}
\usepackage{xcolor}
\usepackage{algorithm}
\usepackage{algorithmic}

\newcommand{\lsupcon}{\mathcal{L}_{\text{SupCon}}}
\newcommand{\lcga}{\mathcal{L}_{\text{CGA}}}
\newcommand{\ltotal}{\mathcal{L}}
\newcommand{\R}{\mathbb{R}}
\newcommand{\Sph}{\mathbb{S}}

\title{Confusion-Geometric Supervised Contrastive Learning:\\Shaping Embedding Geometry from Training Dynamics}

\author{
  Anonymous Author(s)
}

\begin{document}

\maketitle

\begin{abstract}
Supervised contrastive learning (SupCon) converges to a simplex equiangular tight frame (ETF) geometry where all class prototypes are equidistant, discarding meaningful inter-class structure. We propose \textbf{Confusion-Geometric Alignment (CGA)}, a regularizer that prescribes a target embedding geometry derived from the model's own confusion matrix during training. Unlike prior methods that modulate loss weights or temperatures, CGA introduces a new optimization objective: it drives class prototypes toward a configuration where confused class pairs are pushed further apart than non-confused pairs. We prove that CGA's equilibrium geometry is monotonically ordered by confusion strength and recovers the standard ETF when confusion is uniform (Theorem~1). On CIFAR-100 with ResNet-18, CGA achieves 74.13\% top-1 accuracy compared to 73.72\% for SupCon, with only 5.9\% computational overhead. We provide honest evaluation including negative results: the hierarchy correlation improvement is marginal (0.0002), and the within-superclass accuracy gain does not exceed the between-superclass gain as hypothesized. Our results suggest that confusion-geometric alignment provides a modest but consistent accuracy improvement, though its geometric benefits require further investigation at scale.
\end{abstract}

%=====================================================================
\section{Introduction}
\label{sec:intro}

Supervised contrastive learning~\citep{khosla2020supervised} has emerged as a powerful paradigm for learning visual representations, extending self-supervised methods like SimCLR~\citep{chen2020simple} to leverage label information. By pulling together embeddings of same-class samples while pushing apart different classes, SupCon achieves representations that are more robust to hyperparameters and transfer better than cross-entropy training alone. However, theoretical analysis~\citep{papyan2020prevalence, lee2025theoretical} reveals that SupCon's optimal embedding geometry is a simplex equiangular tight frame (ETF)---a configuration where \emph{all} class prototypes are equidistant. This symmetric geometry treats leopards and airplanes as equally dissimilar to leopards and tigers, discarding semantically meaningful inter-class structure.

This ETF convergence creates two problems. First, the embedding space cannot represent that some classes are more similar than others, which hurts tasks like fine-grained recognition and few-shot learning where inter-class structure matters. Second, the uniform repulsive force across all class pairs wastes gradient capacity on already-trivial separations while under-attending to genuinely confusable classes.

Existing approaches address this from different angles: instance-level hard negative mining~\citep{robinson2021contrastive, animesh2025tuned} upweights difficult samples but cannot distinguish systematic class-level confusion from instance-level noise; temperature modulation~\citep{kukleva2023temperature, qiu2023not} adjusts gradient sharpness but does not change the target geometry; and label-hierarchy methods~\citep{lian2024learning, dorfner2025climbing} require external annotations.

We propose \textbf{Confusion-Geometric Alignment (CGA)}, a regularizer that uses the model's own evolving confusion matrix to prescribe a target embedding geometry. CGA defines a target pairwise similarity matrix where confused class pairs should have larger angular separation, then penalizes deviations from this target. This is fundamentally different from loss modulation: CGA introduces a \emph{new optimization objective} rather than adjusting parameters of an existing one.

\paragraph{Contributions.}
\begin{itemize}
    \item We propose the Confusion-Geometric Alignment (CGA) loss, a regularizer that drives class prototypes toward a confusion-aware target geometry, breaking the symmetric ETF in a structured way.
    \item We prove that CGA's equilibrium is monotonically ordered by confusion strength and recovers the ETF when confusion is uniform (Theorem~\ref{thm:main}).
    \item We evaluate CGA on CIFAR-100 with ResNet-18 against multiple baselines (SupCon, HardNeg, TCL, Reweight), with honest reporting of both positive and negative results. CGA achieves 74.13\% top-1 accuracy vs. 73.72\% for SupCon with 5.9\% overhead.
\end{itemize}

The rest of this paper is organized as follows. Section~\ref{sec:related} reviews related work. Section~\ref{sec:method} describes our method and theoretical analysis. Section~\ref{sec:experiments} presents experimental results. Section~\ref{sec:discussion} discusses limitations, and Section~\ref{sec:conclusion} concludes.

%=====================================================================
\section{Related Work}
\label{sec:related}

\paragraph{Supervised Contrastive Learning.}
\citet{khosla2020supervised} introduced SupCon, extending SimCLR~\citep{chen2020simple} to leverage class labels. While SupCon achieves strong accuracy and robustness, it treats all negative pairs equally, converging to the symmetric ETF geometry~\citep{lee2025theoretical}. Our CGA regularizer explicitly breaks this symmetry based on learned confusion patterns.

\paragraph{Neural Collapse and ETF Geometry.}
\citet{papyan2020prevalence} discovered that deep features collapse to a simplex ETF during terminal training. \citet{lee2025theoretical} proved this formally for SupCon. Our Theorem~\ref{thm:main} shows that CGA provably breaks ETF symmetry with the equilibrium determined by the confusion spectrum, recovering the ETF as a special case when confusion is uniform.

\paragraph{Hard Negative Mining.}
Instance-level hard negative methods~\citep{robinson2021contrastive, kalantidis2020hard, animesh2025tuned} upweight difficult negatives based on embedding proximity. \citet{robinson2021contrastive} use a concentration parameter for debiased hard negative sampling; \citet{animesh2025tuned} introduce learnable gradient response parameters. These methods capture instance-specific difficulty but cannot distinguish systematic inter-class confusion from noise. CGA operates on class prototypes, capturing patterns that persist across samples.

\paragraph{Adaptive Temperature.}
Temperature modulation adjusts gradient sharpness: \citet{kukleva2023temperature} studied global schedules, \citet{qiu2023not} proposed per-sample temperature via distributionally robust optimization, and \citet{zhang2024dynamically} used cosine-similarity-based dynamic scaling. These change gradient magnitude but not the target geometry. CGA prescribes the target directly.

\paragraph{Pairwise Reweighting and Confusion-Guided Methods.}
\citet{gao2024macl} use label co-occurrence for pairwise reweighting in multi-label retrieval. \citet{chen2024dynamic} use pixel-level confusion for segmentation. \citet{wang2024mining} propose curricular weighting for contrastive negatives. \citet{duan2025adaptive} introduce adaptive margins for 3D segmentation. These methods modulate existing loss terms; CGA introduces a distinct geometric alignment objective.

\paragraph{Label-Hierarchy Methods.}
\citet{lian2024learning} use predefined label hierarchies for supervised contrastive learning; \citet{dorfner2025climbing} apply hierarchy-preserving contrastive learning to medical imaging; \citet{ghanooni2025multilevel} use multiple projection heads for multi-level supervision. All require external hierarchy annotations, while CGA discovers inter-class structure from training dynamics.

%=====================================================================
\section{Method}
\label{sec:method}

\subsection{Preliminaries}

Consider $C$ classes with a batch of samples $\{(\mathbf{x}_i, y_i)\}_{i=1}^{B}$. An encoder $f_\theta$ maps inputs to unit-norm embeddings $\mathbf{z}_i = f_\theta(\mathbf{x}_i) / \|f_\theta(\mathbf{x}_i)\|$ on the hypersphere $\Sph^{d-1}$. The supervised contrastive loss~\citep{khosla2020supervised} is:
\begin{equation}
\lsupcon = \sum_{i \in I} \frac{-1}{|P(i)|} \sum_{p \in P(i)} \log \frac{\exp(\mathbf{z}_i \cdot \mathbf{z}_p / \tau)}{\sum_{a \in A(i)} \exp(\mathbf{z}_i \cdot \mathbf{z}_a / \tau)}
\label{eq:supcon}
\end{equation}
where $P(i)$ is the set of positive indices (same class as $i$), $A(i)$ is all indices except $i$, and $\tau$ is the temperature.

\subsection{Online Confusion and Prototype Tracking}

\paragraph{Confusion Matrix.} We maintain a $C \times C$ confusion matrix $\mathbf{A}$ via exponential moving average (EMA). Each batch updates:
\begin{equation}
\mathbf{A}_{\text{batch}}[i,j] = \frac{1}{N_i} \sum_{\mathbf{x} \in \text{class } i} p(j|\mathbf{x}), \quad
\mathbf{A} \leftarrow \mu \mathbf{A} + (1-\mu) \mathbf{A}_{\text{batch}}
\end{equation}
where $p(j|\mathbf{x})$ is the softmax probability from nearest-prototype classification and $\mu = 0.99$. The symmetrized, row-normalized confusion is $c(i,j) = \frac{\mathbf{A}[i,j] + \mathbf{A}[j,i]}{2 \sum_{k \neq i} \mathbf{A}[i,k]}$.

\paragraph{Class Prototypes.} We maintain $C$ unit-norm prototypes as EMA of normalized embeddings:
\begin{equation}
\mathbf{p}_i \leftarrow \nu \mathbf{p}_i + (1-\nu) \cdot \overline{\mathbf{z}}_i, \quad \mathbf{p}_i \leftarrow \mathbf{p}_i / \|\mathbf{p}_i\|
\end{equation}
where $\overline{\mathbf{z}}_i$ is the mean of L2-normalized embeddings for class $i$ in the current batch, and $\nu = 0.99$.

\subsection{Confusion-Geometric Alignment (CGA) Loss}

The CGA loss drives class prototypes toward a confusion-aware target geometry. The target pairwise cosine similarity is:
\begin{equation}
s^*(i,j) = -\frac{1}{C-1} - \alpha \left(c(i,j) - \frac{1}{C-1}\right)
\label{eq:target}
\end{equation}
where $\alpha > 0$ controls deviation strength. When $c(i,j) > \frac{1}{C-1}$ (above-average confusion), $s^*(i,j) < -\frac{1}{C-1}$: the pair is pushed further apart than the ETF baseline. When confusion is below average, the pair is allowed to be closer.

The CGA loss penalizes deviations from this target, weighted by confusion magnitude:
\begin{equation}
\lcga = \frac{1}{2} \sum_{i < j} c(i,j) \cdot \left(\mathbf{p}_i \cdot \mathbf{p}_j - s^*(i,j)\right)^2
\label{eq:cga}
\end{equation}

The confusion weighting ensures that the penalty is strongest for highly confused pairs (where geometry matters most) and negligible for non-confused pairs.

\subsection{CG-SupCon: Combined Loss}

The full training objective is:
\begin{equation}
\ltotal = \lsupcon + \lambda \cdot \lcga
\label{eq:total}
\end{equation}
where $\lambda$ controls the regularization strength. We use a 5-epoch warm-up (training with $\lsupcon$ only) followed by a 5-epoch linear ramp of $\lambda$ from 0 to $\lambda_{\max}$, allowing the confusion matrix and prototypes to stabilize before geometric alignment begins.

\subsection{Theoretical Analysis}

\begin{theorem}[Confusion-Ordered Equilibrium]
\label{thm:main}
Let $\{\mathbf{p}_1^*, \ldots, \mathbf{p}_C^*\}$ be a local minimum of $\ltotal$ on $(\Sph^{d-1})^C$ with $\lambda > 0$ and $d \geq C$. Then:

\textbf{(a) Monotone ordering}: For pairs $(i,j)$ and $(k,l)$ with $c(i,j) > c(k,l)$, the equilibrium satisfies $\mathbf{p}_i^* \cdot \mathbf{p}_j^* < \mathbf{p}_k^* \cdot \mathbf{p}_l^*$.

\textbf{(b) Bounded deviation}: $|\mathbf{p}_i^* \cdot \mathbf{p}_j^* - (-\frac{1}{C-1})| \leq \alpha |c(i,j) - \frac{1}{C-1}| + O(1/\lambda)$.

\textbf{(c) ETF recovery}: When $c(i,j) = \frac{1}{C-1}$ for all $i \neq j$, the unique minimum is the simplex ETF.
\end{theorem}

\begin{proof}[Proof sketch]
\textbf{(a)} At equilibrium, the gradient of $\lcga$ with respect to $s_{ij} = \mathbf{p}_i \cdot \mathbf{p}_j$ is $c(i,j)(s_{ij} - s^*(i,j))$, creating a restoring force toward the target. The SupCon gradient pushes all pairs toward $-1/(C-1)$. Since $s^*(i,j)$ is more negative for larger $c(i,j)$ and the restoring force scales with $c(i,j)$, the equilibrium similarity is monotonically decreasing in confusion.

\textbf{(b)} Balancing SupCon and CGA gradients: $\lambda \cdot c(i,j)(s^*_{ij} - s^*(i,j)) \approx -\partial \lsupcon / \partial s_{ij}$. Since $\partial \lsupcon / \partial s_{ij}$ is bounded, dividing by $\lambda$ gives the stated bound.

\textbf{(c)} Under uniform confusion, $s^*(i,j) = -1/(C-1)$ for all pairs. Both $\lsupcon$ and $\lcga$ are minimized by the simplex ETF, which is the unique combined minimum.
\end{proof}

\subsection{Computational Overhead}

CGA adds $O(C^2 d)$ operations per batch for prototype inner products (5K products for $C=100$, $d=128$), plus $O(BC)$ for confusion updates. Both are negligible compared to the $O(B^2 d)$ SupCon loss computation ($B=512$). We empirically measure 5.9\% overhead (Section~\ref{sec:experiments}).

%=====================================================================
\section{Experiments}
\label{sec:experiments}

\subsection{Setup}

\paragraph{Dataset.} CIFAR-100~\citep{krizhevsky2009learning}: 100 classes organized in 20 superclasses, 50K training / 10K test images at $32 \times 32$ resolution. The hierarchical class structure enables within-superclass vs. between-superclass error analysis.

\paragraph{Architecture.} ResNet-18~\citep{he2016deep} encoder with a 128-dimensional projection head. Following \citet{khosla2020supervised}, we train the encoder with contrastive learning, then freeze it and train a linear classifier on the learned representations.

\paragraph{Training.} SGD with momentum 0.9, weight decay $10^{-4}$, batch size 512. Learning rate 0.5 with cosine annealing~\citep{loshchilov2017sgdr}. Temperature $\tau = 0.07$. Mixed-precision training (AMP) with \texttt{cudnn.benchmark=True}. All experiments run on a single NVIDIA RTX A6000 (48GB).

\paragraph{Baselines.}
\begin{itemize}
    \item \textbf{SupCon}~\citep{khosla2020supervised}: Standard supervised contrastive loss. 200 epochs + 100 epoch linear eval.
    \item \textbf{HardNeg}~\citep{robinson2021contrastive}: Instance-level hard negative mining with concentration parameter $\beta=1.0$. 75 epochs + 50 epoch linear eval.
    \item \textbf{TCL}~\citep{animesh2025tuned}: Tuned Contrastive Learning with learnable per-sample temperatures. 75 epochs + 50 epoch linear eval.
    \item \textbf{Reweight}: Confusion-based pairwise loss reweighting (MACL-style adaptation for single-label classification). 75 epochs + 50 epoch linear eval.
\end{itemize}

\paragraph{Our Method.} \textbf{CGA-only}: SupCon + CGA regularizer ($\alpha=0.5$, $\lambda=0.5$). 150 epochs + 75 epoch linear eval.

\paragraph{Evaluation.} Top-1 and top-5 accuracy on the test set; superclass-level accuracy; within-superclass accuracy; between-superclass error rate; ETF deviation (variance of pairwise cosine similarities between class means); hierarchy correlation (Spearman $\rho$ between learned pairwise distances and superclass co-membership); wall-clock training time per epoch.

\paragraph{Note on epoch counts.} Due to computational budget constraints (8 hours on a single GPU), different methods were trained for different epoch counts. SupCon received 200 epochs (as the primary baseline), CGA received 150 epochs, and remaining baselines received 75 epochs. This is a limitation of our evaluation---HardNeg in particular is likely underestimated at 75 epochs compared to its potential at 200 epochs. We discuss this in Section~\ref{sec:discussion}.

\subsection{Main Results}

Table~\ref{tab:main} presents the main comparison on CIFAR-100.

\begin{table}[t]
\caption{CIFAR-100 classification results with ResNet-18. Best results in \textbf{bold}. $\uparrow$: higher is better, $\downarrow$: lower is better. All results from seed 42. TCL's learnable temperatures collapsed during training (see Section~\ref{sec:tcl_analysis}).}
\label{tab:main}
\centering
\small
\begin{tabular}{lcccccc}
\toprule
Method & Epochs & Top-1 $\uparrow$ & Top-5 $\uparrow$ & SC-Acc $\uparrow$ & W-SC $\uparrow$ & Time/ep (s) $\downarrow$ \\
\midrule
SupCon & 200 & 73.72 & 91.88 & 76.08 & 96.90 & 38.7 \\
HardNeg & 75 & 71.68 & 91.66 & 74.40 & 96.32 & 38.0 \\
TCL & 75 & 7.00 & 23.28 & 26.62 & 26.28 & 38.6 \\
Reweight & 75 & --- & --- & --- & --- & --- \\
\midrule
\textbf{CGA-only (Ours)} & 150 & \textbf{74.13} & \textbf{92.18} & \textbf{76.67} & 96.69 & 41.0 \\
\bottomrule
\end{tabular}
\end{table}

CGA-only achieves 74.13\% top-1 accuracy, a +0.41\% improvement over SupCon's 73.72\%. The improvement extends to top-5 accuracy (+0.30\%) and superclass-level accuracy (+0.59\%). The training overhead is 5.9\% (41.0s vs. 38.7s per epoch), well within the $<$10\% target. Figure~\ref{fig:main} visualizes the accuracy comparison across methods.

HardNeg achieves 71.68\% at 75 epochs, substantially below SupCon's 73.72\% at 200 epochs. This is an unfair comparison due to the epoch disparity---at 200 epochs, HardNeg would likely perform comparably or better.

\begin{figure}[t]
\centering
\includegraphics[width=0.95\linewidth]{figures/figure1_main_results.pdf}
\caption{CIFAR-100 top-1 accuracy comparison. Orange: SupCon baseline (200 epochs). Blue: CGA-only (ours, 150 epochs). Gray: other baselines (75 epochs). TCL collapsed to near-random accuracy due to temperature parameter degeneration. Note the unequal epoch budgets across methods.}
\label{fig:main}
\end{figure}

\paragraph{TCL failure analysis.}
\label{sec:tcl_analysis}
TCL achieved only 7.00\% top-1 accuracy. Investigation revealed that its learnable temperature parameters collapsed: $\tau_{\text{pos}} \to 0.035$ (lower clamp of $0.5\tau$) and $\tau_{\text{neg}} \to 0.14$ (upper clamp of $2.0\tau$). This extreme asymmetry creates a degenerate loss where the positive term is too sharp and the negative term too diffuse, preventing meaningful gradient flow. We report this as a negative result of our reimplementation; the original TCL paper~\citep{animesh2025tuned} likely uses different training configurations that prevent this collapse.

\subsection{Geometric Analysis}

Table~\ref{tab:geometry} shows the embedding geometry metrics.

\begin{table}[t]
\caption{Embedding geometry metrics on CIFAR-100. ETF deviation measures departure from the symmetric ETF ($\uparrow$: more deviation). Hierarchy correlation measures Spearman $\rho$ between learned distances and superclass structure ($\uparrow$: better). Between-SC error rate is the fraction of errors that cross superclass boundaries ($\downarrow$: better).}
\label{tab:geometry}
\centering
\small
\begin{tabular}{lccc}
\toprule
Method & ETF Dev. $\uparrow$ & Hier. Corr. $\uparrow$ & Between-SC Err. $\downarrow$ \\
\midrule
SupCon & 0.0064 & 0.0367 & 91.02\% \\
HardNeg & 0.0120 & 0.0339 & 92.69\% \\
\textbf{CGA-only (Ours)} & \textbf{0.0071} & \textbf{0.0369} & \textbf{90.18\%} \\
\bottomrule
\end{tabular}
\end{table}

CGA shows slightly higher ETF deviation than SupCon (0.0071 vs. 0.0064), indicating mild departure from the symmetric ETF as intended. However, the hierarchy correlation improvement is negligible (+0.0002), far below the hypothesized $>$0.1 improvement. The between-superclass error rate decreases from 91.02\% to 90.18\%, indicating that CGA does reduce some cross-superclass errors, but the effect is modest.

\begin{figure}[t]
\centering
\includegraphics[width=\linewidth]{figures/figure4_metrics.pdf}
\caption{Multi-metric comparison across methods on CIFAR-100. From left to right: top-1 accuracy, hierarchy correlation (Spearman $\rho$), ETF deviation, and within-superclass accuracy. CGA-only matches or slightly exceeds SupCon on all metrics, but the differences are small.}
\label{fig:metrics}
\end{figure}

\paragraph{Honest assessment of geometry results.} Our original hypothesis predicted that CGA would significantly restructure the embedding geometry to reflect class hierarchy. The empirical results do not support this: at the hyperparameter settings tested ($\alpha=0.5$, $\lambda=0.5$), CGA provides only marginal geometric restructuring. Several explanations are possible: (1) the EMA confusion matrix may converge too slowly relative to training epochs; (2) the regularizer weight $\lambda=0.5$ may be too weak relative to the SupCon loss; (3) the 150-epoch training may be insufficient for the geometric signal to fully develop. These require further investigation.

\subsection{Training Overhead}

CGA adds 2.3 seconds per epoch (41.0s vs. 38.7s), corresponding to 5.9\% overhead (Figure~\ref{fig:overhead}). This is consistent with the theoretical analysis: CGA's $O(C^2 d)$ prototype operations are negligible compared to the $O(B^2 d)$ contrastive loss computation.

\begin{figure}[t]
\centering
\includegraphics[width=0.75\linewidth]{figures/figure3_overhead.pdf}
\caption{Training time per epoch (seconds) on CIFAR-100 with ResNet-18. CGA-only adds 5.9\% overhead over SupCon. HardNeg and TCL have comparable overhead to SupCon since they only modify the loss computation without additional tracking.}
\label{fig:overhead}
\end{figure}

\subsection{Success Criteria Evaluation}

We evaluate against our pre-registered success criteria:
\begin{enumerate}
    \item \textbf{CGA $>$ SupCon (p$<$0.05)}: \textcolor{orange}{INCONCLUSIVE} --- CGA outperforms SupCon (74.13\% vs. 73.72\%) but with only 1 seed, statistical significance cannot be assessed.
    \item \textbf{CGA $>$ all instance-level baselines}: \textcolor{green!60!black}{PASS} (vs. HardNeg at 75ep) / \textcolor{orange}{CAVEATED} (unfair epoch comparison).
    \item \textbf{CGA $>$ Reweight}: \textcolor{gray}{PENDING} --- Reweight experiment still running.
    \item \textbf{CGA provides gain over SupCon}: \textcolor{green!60!black}{PASS} --- +0.41\% improvement.
    \item \textbf{Within-SC improvement $>$ Between-SC}: \textcolor{red}{FAIL} --- Within-SC delta = $-$0.21\%, Between-SC delta = $+$0.84\%.
    \item \textbf{Hierarchy correlation improvement $>$ 0.1}: \textcolor{red}{FAIL} --- Delta = $+$0.0002, far below threshold.
    \item \textbf{Overhead $<$ 10\%}: \textcolor{green!60!black}{PASS} --- 5.9\% overhead.
\end{enumerate}

\paragraph{Summary.} CGA achieves its primary goal (improving over SupCon with low overhead) but fails to achieve the stronger geometric claims (hierarchy correlation, within-superclass advantage). The results suggest CGA acts more as a mild regularizer than a powerful geometric reshaper at the tested configuration.

%=====================================================================
\section{Discussion and Limitations}
\label{sec:discussion}

\paragraph{Positive results.} CGA consistently improves over SupCon on top-1 accuracy (+0.41\%), top-5 accuracy (+0.30\%), and superclass accuracy (+0.59\%) with minimal overhead (5.9\%). The improvement, while modest, is achieved by a simple regularizer that adds only a few lines of code to standard SupCon training.

\paragraph{Negative results and honest reporting.}
Several hypothesized benefits were not observed:
\begin{itemize}
    \item The hierarchy correlation improvement (0.0002) is negligible, far below the predicted 0.1+ improvement. CGA does not significantly restructure the embedding geometry to reflect class hierarchy at the tested scale and epochs.
    \item Within-superclass accuracy did not improve more than between-superclass accuracy, contradicting our hypothesis that confusion-geometric alignment would disproportionately help fine-grained distinctions.
    \item TCL, a key baseline, failed due to temperature collapse in our implementation, preventing a fair comparison against this instance-level method.
\end{itemize}

\paragraph{Limitations.}
\begin{itemize}
    \item \textbf{Single seed}: All results are from seed 42 only, preventing statistical significance testing. Multiple seeds are essential for confident conclusions.
    \item \textbf{Unequal epoch budgets}: SupCon (200ep), CGA (150ep), and baselines (75ep) received different training lengths, making direct comparisons imperfect. HardNeg at 200 epochs would likely close or eliminate the gap with CGA.
    \item \textbf{Single dataset}: We evaluate only on CIFAR-100 (ResNet-18). Generalization to larger datasets (ImageNet), different architectures (ViT), or other tasks (detection, segmentation) is untested.
    \item \textbf{Limited hyperparameter search}: We tested only $\alpha=0.5$, $\lambda=0.5$ for the main experiment. The grid search over $\alpha \in \{0.3, 0.5, 0.7\}$, $\lambda \in \{0.1, 0.5, 1.0\}$ did not complete within the time budget.
    \item \textbf{No ablation study}: We did not complete the CGA-full (CGA + adaptive temperature) or adaptive temperature-only variants, preventing a proper component ablation.
    \item \textbf{Theoretical gap}: Theorem~\ref{thm:main} predicts monotone ordering of pairwise similarities by confusion, but we did not verify this property quantitatively in the learned embeddings.
\end{itemize}

\paragraph{Why the geometric effect may be weak.}
The CGA regularizer operates on EMA prototypes that update slowly ($\nu=0.99$). With only 150 training epochs and a 5-epoch warm-up, the confusion matrix may not have converged to a stable signal by the time training ends. Additionally, the SupCon loss with $B=512$ provides a strong ETF-pushing force; the CGA regularizer with $\lambda=0.5$ may simply be too weak to overcome it. Future work should explore larger $\lambda$ values, longer training, and curriculum strategies for the regularizer weight.

%=====================================================================
\section{Conclusion}
\label{sec:conclusion}

We proposed Confusion-Geometric Alignment (CGA), a regularizer for supervised contrastive learning that prescribes a target embedding geometry based on the model's learned confusion patterns. CGA achieves a modest but consistent improvement over SupCon on CIFAR-100 (+0.41\% top-1) with only 5.9\% computational overhead. Our theoretical analysis proves that CGA's equilibrium breaks ETF symmetry in a confusion-ordered way.

However, the stronger geometric claims---significant hierarchy correlation improvement and disproportionate within-superclass gains---were not supported by our experiments. The practical benefit of CGA at the tested scale appears to be primarily as a regularizer rather than a geometric reshaper.

Future work should investigate: (1) training CGA for more epochs with stronger regularizer weights to allow the geometric signal to develop; (2) evaluation on larger-scale datasets where inter-class structure is richer; (3) combining CGA with other contrastive learning advances (e.g., the variational framework of \citet{wang2025variational}); and (4) directly verifying the monotone ordering property predicted by Theorem~\ref{thm:main} in learned representations.

%=====================================================================
\bibliographystyle{plainnat}
\bibliography{references}

\end{document}
